\documentclass[final,hyperref={pdfpagelabels=false}]{beamer}
\mode<presentation>
  {
  %  \usetheme{Berlin}
  \usetheme{Dreuw}
  }
  \usepackage{times}
  \usepackage{amsmath,amsthm, amssymb, latexsym, tikz}
\usepackage{graphicx}

\usepackage{tabulary}

\usepackage[export]{adjustbox}


\usepackage{multicol}
\usetikzlibrary{arrows,shapes}
  
  \boldmath
  \usepackage[english]{babel}
  \usepackage[latin1]{inputenc}
  \usepackage[orientation=portrait,size=a0,scale=1.4,debug]{beamerposter}

  %%%%%%%%%%%%%%%%%%%%%%%%%%%%%%%%%%%%%%%%%%%%%%%%%%%%%%%%%%%%%%%%%%%%%%%%%%%%%%%%%5
  \graphicspath{{figures/}}
  \title{Making research reproducible \\using literate programming}
  \author{Sebastian Sauer, Sandra S\"ulzenbr\"uck, Yvonne Ferreira, and Christoph Kurz}
  \institute{FOM University of Applied Sciences, Helmholtz Zentrum M\"unchen}

  \date{March 2014}


  %%%%%%%%%%%%%%%%%%%%%%%%%%%%%%%%%%%%%%%%%%%%%%%%%%%%%%%%%%%%%%%%%%%%%%%%%%%%%%%%%5
  \begin{document}
  \begin{frame}{} 
    
      \begin{columns}[t]
      
      \begin{column}{.48\linewidth}
      
        \begin{block}{Reproducibility: What and why?}
      There is an increasing concern about the reliability of research results \cite{Peng2015}.
    Precisely, it has been found that many published results cannot be replicated \cite{OpenScienceCollaboration2015}.
    In parts, this can be due to the fact that it is often hardly possible for independent researchers to
    confirm the results of a research paper. Thus, making research more reproducible seems a pressing need  \cite{Peng2015}.
    Here, we present some ``recipe'' for making (your) research paper more reproducible using literate programming.
    Literate programming refers to weaving programming code (ie., statistical calculations) with the paper text.        
    
    \end{block}
      
      
      
            \begin{block}{Workflow for reproducible paper writing}
         
      \begin{minipage}[t]{0.45\textwidth}
      %Workflow, in general
      %\vspace{2cm}
      \tikzstyle{format} = [rectangle, draw, text centered, rounded corners]
      \tikzstyle{medium} = [rectangle, draw, thin]
      \tikzstyle{line} = [draw, -latex']
      \begin{tikzpicture}[node distance=7cm, auto]
          % We need to set a bounding box first. Otherwise the diagram
          % will change position for each frame.
          %\path[use as bounding box] (-5,0) rectangle (10,10);
        
        \node [format] (idea) {idea};
        \node [format, below of = idea] (notes) {notes};
        \node [format, below of = notes, align = center] (manuscript) {unformated\\manuscript};
        \node [format, below of = manuscript,  align = center] (paper) {typesetted\\paper};
        \node [medium, left of = manuscript, xshift = -3cm] (R)  {calculations};
        \node [medium, above left  of = manuscript, yshift = 1cm, align = center] (vc) {version\\control};
        \node [medium, below left of = manuscript,  xshift = -3cm, yshift = -3cm]  (collab) {collaborators};
        \path [line] (idea) -- node [align=center] {\emph{think,}\\\emph{scribble}} (notes) ;
        \path [line] (notes) --  node [align=center] {\emph{hack}\\\emph{text}} (manuscript);
        \path [line] (manuscript) -- node {\emph{format}}(paper);     
        \path [line, below=7cm] (R) --  node [align=center] {\emph{calc}\\\emph{data}}(manuscript);  
        \path [line, left] (vc) -- node {\emph{track changes}} (manuscript);
        \path [line] (collab) -- node [align=center, near start] {\emph{give}\\\emph{input}} (manuscript);
      \end{tikzpicture}
     \end{minipage}	
     %
     % now with software tools
     \begin{minipage}[t]{0.50\textwidth}
     \begin{tikzpicture}[node distance=7cm, auto]
          % We need to set a bounding box first. Otherwise the diagram
          % will change position for each frame.
          %\path[use as bounding box] (-5,0) rectangle (10,10);
          \tikzstyle{format} = [rectangle, draw, text centered, rounded corners]
      \tikzstyle{medium} = [rectangle, draw, thin]
      \tikzstyle{line} = [draw, -latex']
        \node [format] (idea) {idea};
        \node [format, below of = idea,  align = center] (notes) {Piece\\of Paper};
        \node [format, below of = notes, align = center] (manuscript) {ugly\\.txt file};
        \node [format, below of = manuscript,  align = center] (paper) {polished\\.pdf file};
        \node [medium, left of = manuscript, xshift = -3cm] (R)  {R-code};
        \node [medium, above left  of = manuscript, yshift = 1cm, align = left] (vc) {git};
        \node [medium, below left of = manuscript,  xshift = -3cm, yshift = -3cm]  (collab) {collaborators};
        \path [line] (idea) -- node [align=center]{\emph{think,}\\\emph{scribble} }(notes) ;
        \path [line] (notes) --  node [align=center]{\emph{hack}\\\emph{.md file}}  (manuscript);
        \path [line] (manuscript) -- node [align=center]{\emph{apply}\\\emph{.tex style}}(paper);     
       \path [line, below=7cm] (R) --  node [align=center] {\emph{weave}\\\emph{in}}(manuscript);  
        \path [line, left] (vc) -- node {\emph{track changes}} (manuscript);
        \path [line] (collab) -- node [align=center, near start] {\emph{give}\\\emph{input}} (manuscript);
      \end{tikzpicture}
     \end{minipage}	
     \\
    The \emph{diagram} shows the workflow \emph{in theory}. The \emph{right} diagram exemplifies useful tools for each step.
     
    \end{block}





        \begin{block}{Requirements to reproducible paper writing}
        Some (subjective) criteria to tools helping writing scientific papers reproducibly.
          \begin{itemize}
          \item \textbf{simple:} The tool should be easy to learn and use (flat learning curve). Life is short.
          \item \textbf{beautiful:} The tool should set the text in an eye pleasing way, eg., no holes in block aligned text.
	   \item \textbf{plain:} The tool should work with plain text files, as they will probably readable in (distant) future. In addition, text file are more compatible with other tools, such as version control. 
         \item \textbf{versionized:} The tool should should support versionizing (tracking changes), including changes from some more or less helpful collaborators.
         \item \textbf{citable:} The tool should be able to manage (scholarly) citations.
         \item \textbf{flexible:} The tool should allow for flexible layout, ie., many possibilities for rendering the final, formatted output. This is often in outright contradiction to easyness.
	
	
       \end{itemize}
        \end{block}



   
     \begin{block}{Which tools make the researcher happy?}
    
\begin{tabular}{rllll}
  \hline
 & Word & Latex & Markdown & WebApps \\ 
  \hline
simple &  $\bullet\bullet\bullet$ & $\bullet$ & $\bullet\bullet$ & $\bullet\bullet$ \\ 
  beautiful & $\bullet$ & $\bullet\bullet\bullet$ & $\bullet\bullet\bullet$ &  $\bullet\bullet$ \\ 
  plain & $\bullet$ & $\bullet\bullet\bullet$ & $\bullet\bullet\bullet$ & $\bullet$ \\ 
  %readable & $\bullet$ &  $\bullet\bullet$ & $\bullet\bullet\bullet$ & $\bullet\bullet\bullet$ \\ 
  versionized & $\bullet$ & $\bullet\bullet\bullet$ & $\bullet\bullet\bullet$ & $\bullet\bullet\bullet$ \\ 
  citable & $\bullet\bullet\bullet$ & $\bullet\bullet\bullet$ & $\bullet\bullet\bullet$ & $\bullet\bullet\bullet$ \\ 
  flexible & $\bullet\bullet\bullet$ &  $\bullet\bullet$ & $\bullet$ &  $\bullet\bullet$ \\ 
   \hline
   $\Sigma$ & 13 & 17 & 18 & 16\\
   \hline
\end{tabular}
\bigskip

This table compares some (subjective) criteria of what a researcher needs for writing a paper in a reproducible and efficient way. Note: \emph{Word} refers not only to MS Word, but to similar WYSIWYG text processors as well. \emph{WebApps} refer to scholarly writing tools such as \emph{Authorea}. \emph{Markdown} refers to the \emph{pandoc} dialect and extensions of markdown.
   
        \end{block}
 
      
      \end{column}
      
            
      
      \begin{column}{.48\linewidth} 
          
      
  \begin{block}{Software tools for reproducible writing}
    
\begin{tabulary}{1.0\textwidth}{CLL}
  \hline
 Logo  & Name & Description  \\ 
 \hline
\includegraphics[valign=T,scale=.3]{logo/R_logo} & \textbf{R }& Environment for statistical calculations and programming\\
 \includegraphics[valign=T,scale=.1]{logo/RStudio_logo} &  \textbf{R-Studio} & IDE for R including features for version control (eg git), and literal programming (Knitr, Pandoc, Markdown)\\
 \includegraphics[valign=T, scale=.5]{logo/knitr_logo} &  \textbf{Knitr} & R-package to ``knit'' (or ``weave'') R-code into plain text as a way of functional programming\\
 \includegraphics[valign=T, scale=.3]{logo/git_logo} & \textbf{Git} & version control tool, handles multiples collaborators, tracks changes in text files\\
 \texttt{PANDOC} &  \textbf{Pandoc} & ``swiss army knife'' for converting markup languages, eg. Markdown $\rightarrow$ PDF\\
 %\includegraphics[valign=T, scale=.3]{logo/markdown_logo} & \textbf{Markdown} & Similar aims as Latex, but you learn it in 5 minutes. See: \textbf{boldface} in Markdown\\  
 \hline
  \end{tabulary}
  
   \end{block}



        \begin{block}{code bits'n'bites}
        
          The main commands are:
          \begin{itemize}
             \item Weave code in text using \textbf{knitr}, giving a R-Markdown (.rmd) (or R-Latex, .rnw) file: ``Mean reaction time was \texttt{`round(mean(rt), 2)'} ms''. $\rightarrow$ ``Mean reaction time was 433.30 ms''.
             \item Knit R-Text-Mixup to pure text plus markup using R-package \textbf{knitr:}
             \newline \texttt{knitr(source.rmd} 
             \item Convert Markdown to Latex-PDF with \textbf{Pandoc:} 
             \newline \texttt{pandoc source.md --output paper.pdf ...} 
      \end{itemize}
          
        \end{block}
        
        
        
        \begin{block}{Conclusion}
            \begin{itemize}
             \item A somewhat ideal combinations of tools for efficiently and reproducibly writing papers is: R+RStudio+knitr+markdown+pandoc+git.    
             \item Writing in \texttt{Markdown} is nice, because little markup clutters the view. Pandoc them compiles the markdown to latex, applying some tex-template of your flavor.
             \item However,  \textt{Pandoc} currently does not provide (easily) provides the flexibility to adjust every detail one may think of, in addition changing the default templates is not yet stable (enough for me). 
             \item Looking forward to the time where you can write in R-Markdown and compile it to Tex easily with all kinds of style-templates.
        \end{block}     

   
   
         \end{column}
    \end{columns}
  \end{frame}
\end{document}


%%%%%%%%%%%%%%%%%%%%%%%%%%%%%%%%%%%%%%%%%%%%%%%%%%%%%%%%%%%%%%%%%%%%%%%%%%%%%%%%%%%%%%%%%%%%%%%%%%%%
%%% Local Variables: 
%%% mode: latex
%%% TeX-PDF-mode: t
%%% End:
